% ===========================================================
% -----------------------------------------------------------
% Dedicatória, Agradecimentos e Epígrafe
% -----------------------------------------------------------
% ===========================================================

% -----------------------------------------------------------
% Dedicatória
% -----------------------------------------------------------
%\begin{dedicatoria}
 %  \vspace*{\fill}
  % \centering
   %\noindent
%   \textit{ Este trabalho é dedicado às crianças adultas que,\\
 %  quando pequenas, sonharam em se tornar cientistas.} \vspace*{\fill}
%\end{dedicatoria}

% -----------------------------------------------------------
%  Agradecimentos
% -----------------------------------------------------------
\begin{agradecimentos}
Dedico minha imensa gratidão a todos que, de alguma forma, participaram e foram cruciais para o meu crescimento e para a concretização desta etapa. Um agradecimento especial à minha família, que representou o meu alicerce, fornecendo apoio incondicional, estímulo e compreensão durante toda a jornada acadêmica.

Expresso meu reconhecimento também aos professores que me transmitiram seus conhecimentos e, de forma especial, ao meu orientador, Professor Fábio Lopes Licht, por ter aceitado minha proposta e dedicado seu tempo para me guiar na elaboração deste trabalho.


\end{agradecimentos}


% -----------------------------------------------------------
% Epígrafe
% -----------------------------------------------------------
% Epígrafe é um título ou frase que serve de tema ou introdução de assunto. 
% O termo tem origem no Grego "epigrafhé" que significa "inscrição", "título".
%
% Em uma obra literária, epígrafe é uma curta citação colocada em uma página no 
% início da obra ou em destaque na abertura de um capítulo. Funciona como um 
% resumo do que se vai ler em seguida.
%
% Ref: http://www.significados.com.br/epigrafe/
% -----------------------------------------------------------
%\begin{epigrafe}
%    \vspace*{\fill}
%	\begin{flushright}
	%	\textit{``Não vos amoldeis às estruturas deste mundo, \\
	%		mas transformai-vos pela renovação da mente, \\
		%	a fim de distinguir qual é a vontade de Deus: \\
		%	o que é bom, o que Lhe é agradável, o que é perfeito.\\
		%	(Bíblia Sagrada, Romanos 12, 2)}
	%	\end{flushright}
%	\end{epigrafe}

